\begin{abstract}
We prove Gaussian scaling limits for the two-dimensional discrete Gaussian
height model throughout its rough phase, including the physical roughening
threshold, for every finite square-symmetric interaction set containing the
nearest neighbors. After the natural lattice normalization, the critical
effective temperature has the universal value \(8\pi\).
For the ordinary nearest-neighbor Villain and XY models at sufficiently low
fixed temperatures, we also prove that their Green-function-normalized
spin fields converge to the imaginary exponential of a Dirichlet Gaussian
free field.
\end{abstract}

\section{Introduction and statements}
\label{main:introduction}

Two-dimensional height and planar spin models have phases in which
fluctuations remain visible at every scale.  The expected large-scale
description is a Gaussian free field for heights and an imaginary
exponential of a Gaussian free field for planar spins.  In each case,
an effective parameter derived from the lattice interaction governs the
Gaussian covariance.  The Berezinskii--Kosterlitz--Thouless (BKT) picture
predicts a universal critical value after the model's natural
normalization, although the parameter varies throughout the critical
phase.  Establishing a field
limit in a perturbative regime and identifying that endpoint for the
original lattice model are separate problems.  This paper addresses
both for discrete Gaussian heights, and proves the full spin-field
limit for the ordinary XY and Villain models at sufficiently low fixed
temperature.

\subsection{Height model and its physical threshold}

An \emph{admissible interaction} is a finite set
$J\subset\Z^2\setminus\{0\}$ invariant under every rotation and reflection
of the square and containing
$J_{\nn}=\{(1,0),(-1,0),(0,1),(0,-1)\}$.  Define
\begin{equation}\label{main:laplacian}
 (\Delta_Ju)(x)=\frac1{|J|}\sum_{y\in J}\bigl(u(x+y)-u(x)\bigr),
 \qquad
 v_J^2=\frac1{2|J|}\sum_{y\in J}y_1^2.
\end{equation}
The height variable takes values in $2\pi\Z$ and $\beta>0$ is its
\emph{temperature}.  For $Q_R=[-R,R]^2\cap\Z^2$ let
$\mu^0_{J,\beta,R}$ be the law with heights zero outside $Q_R$ and weight
\begin{equation}\label{main:zero-law}
 \exp\left\{-\frac1{4\beta|J|}
    \sum_{x\in\Z^2}\sum_{y\in J}
        (\sigma_{x+y}-\sigma_x)^2\right\}.
\end{equation}
Both orientations of each edge occur in the sum.  The threshold is
\begin{equation}\label{main:physical-threshold}
 \bc(J)=\inf\left\{\beta>0:
       \sup_{R\ge1}\E_{\mu^0_{J,\beta,R}}\sigma_0^2=\infty\right\}.
\end{equation}

For $n=L^N$ write $\Lambda_N=(\Z/n\Z)^2$.  The torus law is
\begin{equation}\label{main:torus-law}
 \Prob_{J,\beta,N}(\sigma)=\frac1Z
       \exp\left\{-\frac{(\sigma,-\Delta_J\sigma)}{2\beta}\right\},
 \qquad \sigma\in(2\pi\Z)^{\Lambda_N},\quad \sigma_0=0,
\end{equation}
where $(u,v)=\sum_xu_xv_x$ and steps are taken modulo $n$.
Pinning chooses a representative of the global height shift.
Set $\bar\sigma=n^{-2}\sum_x\sigma_x$ and
\begin{equation}\label{main:height-field}
 H_N=n^{-2}\sum_{x\in\Lambda_N}(\sigma_x-\bar\sigma)\delta_{x/n}.
\end{equation}
On the unit torus $\T^2=\R^2/\Z^2$, let $\Phi_{\T}$ be the mean-zero
real Gaussian free field with covariance $(-\Delta_{\T})^{-1}$,
where $\Delta_{\T}$ is the Euclidean Laplacian.  Pairings of continuum
functions use Lebesgue measure.

\begin{theorem}[Height universality through the roughening threshold]
\label{main:heights}
For every admissible $J$, the threshold in
Equation~\eqref{main:physical-threshold} satisfies $0<\bc(J)<\infty$.
The supremum there is finite when $0<\beta<\bc(J)$ and infinite when
$\beta\ge\bc(J)$, including equality.

There are an integer $L=L(J)\ge2$ and a unique positive function
$\beff(J,\beta)$, defined for $\beta\ge\bc(J)$, such that along the
full sequence $N\to\infty$
\begin{equation}\label{main:height-limit}
 H_N\ \Longrightarrow\
       \sqrt{\frac{\beff(J,\beta)}{v_J^2}}\,\Phi_{\T}
       \quad\hbox{in }H^{-3}(\T^2).
\end{equation}
For every smooth real mean-zero $f$,
\begin{equation}\label{main:height-laplace}
 \lim_{N\to\infty}\log\E e^{H_N(f)}
 =\frac{\beff(J,\beta)}{2v_J^2}
           (f,(-\Delta_{\T})^{-1}f).
\end{equation}
Furthermore,
\begin{align}
 \beff(J,\bc(J))&=8\pi v_J^2,\label{main:height-endpoint}\\
 \lim_{\beta\downarrow\bc(J)}
 \frac{\beff(J,\beta)-\beff(J,\bc(J))}{\beta-\bc(J)}
    &=+\infty.\label{main:height-slope}
\end{align}
The function $\beff(J,\cdot)$ is differentiable on $(\bc(J),\infty)$,
and for some $c_J>0$,
\begin{equation}\label{main:height-high-temperature}
 \beff(J,\beta)=\beta+O_J(e^{-c_J\beta})\qquad(\beta\to\infty).
\end{equation}
For $J_\rho=\{x\in\Z^2\setminus\{0\}:|x|_\infty\le\rho\}$,
\begin{equation}\label{main:height-spread}
 \frac{\bc(J_\rho)}{8\pi v_{J_\rho}^2}\longrightarrow1
      \qquad(\rho\to\infty).
\end{equation}
In particular $v_{J_{\nn}}^2=1/4$ and the nearest-neighbor effective
temperature at roughening is $2\pi$.
\end{theorem}

\paragraph{Other torus sizes and the localized phase.}
The internal transfer theorem gives a stronger torus statement than
the geometric-sequence assertion in Theorem~\ref{main:heights}.
Fix an admissible $J$ and any $k=\beta/v_J^2>0$, and let $a(k)\ge0$ be
the structural covariance coefficient constructed in
Section~\ref{st:section}.  Replace $n=L^N$ in the torus law and
mean-subtracted field above by any increasing sequence of integer
side lengths.  Theorem~\ref{an:dirichlet-torus} proves convergence in
$H^{-3}(\T^2)$ to the mean-zero GFF of covariance
$a(k)(-\Delta_{\T})^{-1}$, together with convergence of the Laplace
transforms of smooth mean-zero tests.  When $a(k)=0$, the limit is
the zero distribution.  On the positive branch,
$\beff(J,\beta)=v_J^2a(k)$, so the same coefficient governs every
such sequence.  Corollary~\ref{an:physical-criterion} identifies this
branch with physical roughness; below $\bc(J)$, the zero limit
expresses the disappearance of macroscopic height fluctuations at
this normalization.  The proof of the transfer theorem works directly
with arbitrary lattice sizes.  On its zero branch the covariance
upper bound alone suffices, while the positive branch also uses the
Dirichlet lower bound.  This stronger result is established before
the endpoint analysis and retains the mean subtraction in
Equation~\eqref{main:height-field}.

\subsection{Ordinary planar spin models and the Wick convention}
\label{main:spin-models}

Here $b>0$ is \emph{inverse temperature}.  Set $D=(-1,1)^2$ and
$D_n=[-1,1]^2\cap n^{-1}\Z^2$, for every integer $n\ge1$.
Boundary vertices have a coordinate equal to $-1$ or $1$ and their
angles are fixed to zero.  Each interior angle lies in
$\R/(2\pi\Z)$.  Write $E_n$ for the unordered nearest-neighbor pairs
in $D_n$, including edges incident to its fixed boundary.
With respect to product Haar measure on the interior angles, the
two probability densities are proportional to
\begin{align}
 \text{XY:}\quad &\prod_{\{x,y\}\in E_n}e^{b\cos(\theta_x-\theta_y)},
       \label{main:xy-law}\\
 \text{Villain:}\quad &\prod_{\{x,y\}\in E_n}
       \sum_{m\in\Z}e^{-b(\theta_x-\theta_y+2\pi m)^2/2}.
       \label{main:villain-law}
\end{align}
Every edge occurs once.  Let $\mathcal L_n$ be the unscaled graph
Laplacian on the interior, with zero boundary values:
\begin{equation}\label{main:spin-green}
 (\mathcal L_nu)(x)=\sum_{y\in D_n:\,|y-x|=1/n}(u(x)-u(y)),
 \qquad G_n=\mathcal L_n^{-1}.
\end{equation}
In particular no factor $n^2$ is included in $\mathcal L_n$.

Let $\Phi_D$ be the zero-Dirichlet GFF with covariance
$(-\Delta_D)^{-1}$; its local covariance singularity is
$(2\pi)^{-1}\log|z-w|^{-1}$.
For $K>1/(4\pi)$ define the complex random distribution
\begin{equation}\label{main:chaos}
 V_{K,D}=\lim_{\epsilon\downarrow0}
     \exp\left\{\frac{i\Phi_{D,\epsilon}}{\sqrt K}
                  +\frac{\E\Phi_{D,\epsilon}^{\,2}}{2K}\right\},
\end{equation}
where $\Phi_{D,\epsilon}$ is the circle average and the limit is tested
against smooth functions of compact support in $D$.  The proof below
also constructs this limit; see \citet{JunnilaSaksmanWebb2020}
for the general imaginary-chaos theory.  The full variance occurs in
Equation~\eqref{main:chaos}, including the finite spatial part of the
Dirichlet Green function.  Consequently
$\E V_{K,D}(f)=\int_D f$.

\begin{theorem}[Full low-temperature spin-field limits]
\label{main:spins}
For each $M\in\{\mathrm{Villain},\mathrm{XY}\}$ there are a finite
$b_0(M)>0$ and functions $K_M(b)>1/(4\pi)$ and $A_M(b)>0$,
defined for all $b\ge b_0(M)$, with $K_M(b)/b\to1$ as $b\to\infty$,
such that
\begin{equation}\label{main:spin-limit}
 S_{M,b,n}:=n^{-2}\sum_{x\in D_n^\circ}
       A_M(b)e^{G_n(x,x)/(2K_M(b))}e^{i\theta_x}\delta_x
       \ \Longrightarrow\ V_{K_M(b),D}
\end{equation}
along the full sequence of integers $n\to\infty$, with tightness and
convergence in $H^{-3}_{\mathrm{loc}}(D)$.
For every finite collection of complex smooth compactly supported test
functions, their smeared fields converge jointly, including real and
imaginary parts.  The coefficients $K_M(b)$ and $A_M(b)$ are independent
of the test functions, the insertion locations, and $n$.
\end{theorem}

\paragraph{All mixed smeared moments.}
For each model, the proof gives more than joint weak convergence.
Fix any $b\ge b_0(M)$ in the asserted range and any finite list
$f_1,\ldots,f_r\in C_c^\infty(D;\C)$, with the same zero boundary
angles and normalization as in Theorem~\ref{main:spins}.  Every mixed
moment of the variables $S_{M,b,n}(f_j)$ and their complex conjugates
converges to the corresponding moment of $V_{K_M(b),D}(f_j)$ and its
conjugate.  Equivalently, expectations of every fixed polynomial in
these variables and their conjugates converge.  Thus the result also
controls integrated multipoint correlations of arbitrary fixed order.
The lower bound $b_0(M)$ is chosen once, independently of that order;
the temperature remains fixed as all integers $n\to\infty$.

Weak convergence alone does not control these unbounded polynomial
observables.  Theorem~\ref{vi:field-transfer}, applied in both
Theorems~\ref{vi:villain-spin} and~\ref{xyf:spin}, supplies the additional
control.  Away from collisions, the separated signed multipoint limits
pass to integrals by Riemann summation.  Near collisions, positivity,
an integrable two-point bound, and Lee--Yang moment bounds make the
omitted contribution vanish, including repeated lattice sites.
The uniform bounds on higher smear moments give uniform integrability
of each fixed polynomial.  The Lee--Yang control of their growth with
moment order also makes the joint limiting law determined by its
moments.  This is the step that connects pointwise correlation information
to both full field convergence and convergence of its smeared moments.

\begin{theorem}[Villain coefficient bridge]\label{main:bridge}
The choices in Theorem~\ref{main:spins} can be made so that
$\pi^2b\ge\bc(J_{\nn})$ throughout the asserted Villain range and
\begin{equation}\label{main:bridge-equality}
 K_{\mathrm{Villain}}(b)
       =\frac{\beff(J_{\nn},\pi^2b)}{\pi^2}.
\end{equation}
\end{theorem}

The spin theorems concern sufficiently large fixed inverse temperatures.
They assert neither a spin-field limit at the BKT endpoint nor equality
of the two models' effective coefficients at the same bare $b$.
The bare conversion in Equation~\eqref{main:bridge-equality} is useful
for checking constants: Fourier duality gives integer heights $h$ with
weight $\exp\{-\sum_e(\nabla_eh)^2/(2b)\}$; setting $\sigma=2\pi h$
in Equation~\eqref{main:torus-law}, with
$-\Delta_{J_{\nn}}=\mathcal L/4$, gives $\beta=\pi^2b$.
The effective conversion requires the source and disorder calculations
in Section~\ref{vi:section}.

In these conventions, the nearest-neighbor height endpoint
$\beff(J_{\nn},\bc)=2\pi$ gives
$\beff(J_{\nn},\bc)/\pi^2=2/\pi$, the dimensionless stiffness value
in the Nelson--Kosterlitz prediction \citep{NelsonKosterlitz1977}.
This arithmetic identifies the normalization correspondence.  The
height theorem reaches the physical roughening threshold, whereas
Theorem~\ref{main:bridge} identifies the spin coefficient only in its
stated sufficiently-low-temperature range.  Extending that bridge and
the spin-field limit to a spin transition would require additional
results.

\subsection{Normalizations and order of limits}

Table~\ref{main:normalizations} collects the conversions used throughout
the proof.  Table~\ref{main:limit-order} specifies the theorem-level
limiting procedures; in particular the spin inverse temperature is
held fixed while the lattice spacing tends to zero.  The auxiliary
choices for annular transfer, physical entry, and endpoint response
are given with their constructions in Sections~\ref{an:section},
\ref{en:section}, and~\ref{ep:section}, respectively.

\begin{table}[htbp]
\centering\small
\begin{tabular}{@{}p{.23\textwidth}p{.70\textwidth}@{}}
\toprule
Quantity & Convention and conversion \\
\midrule
Height law & $\sigma\in2\pi\Z$; temperature $\beta$;
precision $(-\Delta_J)/\beta=A_0/k$, where
$A_0=-\Delta_J/v_J^2$ and $k=\beta/v_J^2$;
$k_c=\bc(J)/v_J^2$. \\[3pt]
Height limit & Covariance $a(k)(-\Delta)^{-1}$ with
$\beff=v_J^2a(k)$; $a(k_c)=8\pi$ and
$\beff(J,\bc)=8\pi v_J^2$. \\[3pt]
Nearest neighbors & $v_{J_{\nn}}^2=1/4$,
$-\Delta_{J_{\nn}}=\mathcal L/4$, and the endpoint
$\beff(J_{\nn},\bc)=2\pi$. \\[3pt]
Spin law & Inverse temperature $b$; each unordered edge counted once;
boundary angle zero; $G_n=\mathcal L_n^{-1}$ with no $n^2$ factor. \\[3pt]
Wick field & $\operatorname{Cov}(\Phi_D)=(-\Delta_D)^{-1}$,
logarithmic coefficient $1/(2\pi)$; exponent $i\Phi_D/\sqrt K$;
normalization uses the full variance, so $\E V_{K,D}=1$. \\[3pt]
Villain bridge & Dual $h\in\Z$, energy $\sum_e(\nabla_e h)^2/(2b)$;
$\sigma=2\pi h$, $\beta=\pi^2b$, and
$K_{\rm Villain}(b)=\beff(J_{\nn},\pi^2b)/\pi^2$. \\[3pt]
XY coefficient & $K_{\rm XY}(b)$ is obtained from the XY construction
in Section~\ref{xyf:section}; $K_M(b)/b\to1$ for each model. \\
\bottomrule
\end{tabular}
\caption{Normalization dictionary.  Temperature and inverse temperature
have different roles in the two lattice descriptions.}
\label{main:normalizations}
\end{table}

\begin{table}[htbp]
\centering\small
\begin{tabular}{@{}p{.23\textwidth}p{.70\textwidth}@{}}
\toprule
Argument & Quantifiers and order \\
\midrule
Height field & Fix $J$, then $L=L(J)$ and $\beta\ge\bc(J)$;
let $N\to\infty$ along every $n=L^N$.  The pinned representative is
mean-subtracted before testing. \\[3pt]
Height endpoint & Define $\beff$ by the field limit first;
then let $\beta\downarrow\bc(J)$ from above.  The spread-out
limit $\rho\to\infty$ concerns the thresholds of the distinct $J_\rho$. \\[3pt]
Spin field & Fix $M$ and any $b\ge b_0(M)$; let every integer
$n\to\infty$.  Circle regularization $\epsilon\downarrow0$ defines
the continuum target locally inside $D$. \\
\bottomrule
\end{tabular}
\caption{Limits in the main results.  The finite Gaussian comparison
statements themselves require no thermodynamic limit.}
\label{main:limit-order}
\end{table}

\subsection{Context and earlier work}

The physical picture originates with Berezinskii and
Kosterlitz--Thouless.  In a low-temperature quadratic approximation,
Berezinskii derived Gaussian multipoint formulas and algebraic spin
correlations with a temperature-dependent exponent
\citep{Berezinskii1971}.  Kosterlitz and Thouless separated smooth angular
fluctuations, or spin waves, from vortices in an approximate description
of the planar model \citep{KosterlitzThouless1973}.  A vortex is a defect
around which the spin angle has nonzero integer winding.  In their BKT
scenario, oppositely charged vortices remain paired at low temperature,
while spin waves produce a critical phase with a continuously varying
decay exponent.  Heating allows vortex pairs to separate; large-distance
screening by the unbound vortices destroys the algebraic behavior.

Kosterlitz's recursion analysis follows the coupled evolution of vortex
activity and stiffness under changes of scale \citep{Kosterlitz1974}.
Here activity measures the weight of vortex defects, and stiffness
measures the response to a slowly varying twist of the angle field.
Bound pairs can change this response, so the coefficient visible at
large distances is renormalized.  Nelson and Kosterlitz predicted the
universal limit $2/\pi$ for the dimensionless renormalized stiffness
from the low-temperature side, and expected it to vanish above the
transition \citep{NelsonKosterlitz1977}.

Fr\"ohlich and Spencer established the transition for several
two-dimensional lattice models, including the planar rotator and
Villain models \citep{FrohlichSpencer1981}.  Their multiscale expansion
organizes neutral collections of charges and controls their cumulative
effect.  A detailed modern account for the discrete Gaussian and Villain
cases is given by Kharash and Peled \citep{KharashPeled2017}.
These phase bounds leave the limiting field and its effective coefficient
unidentified.

Dimock and Hurd controlled the infrared renormalisation flow for the
sine-Gordon model with a fixed ultraviolet cutoff, for each fixed
coefficient above $8\pi$ in their normalization and sufficiently small
activity, with estimates uniform in torus volume
\citep{DimockHurd2000}.  Falco constructed a critical renormalisation
trajectory at small activity for the lattice Coulomb gas and obtained
critical two-point asymptotics for fractional charges, including
logarithmic corrections \citep{Falco2012,Falco2013}.  Using this
approach at the physical threshold of a fixed lattice model requires a
controlled passage from microscopic data to a dilute effective system.
This matching problem already appears in the discussion of
\citet{KosterlitzThouless1973}.

For discrete Gaussian heights, Bauerschmidt, Park, and Rodriguez
smooth the microscopic law to exponentially small activity and prove
the macroscopic torus Gaussian limit at sufficiently high temperature
\citep{BPR2024I}.  They also developed infinite-volume and mesoscopic
scaling limits in \citet{BPR2024II}.  Park subsequently
proved high-temperature central limit theorems for microscopic
multipoint observables \citep{Park2025Multipoint}.
Conjecture~1.3 of \citet{BPR2024I} predicts the torus Gaussian limit
through a critical temperature for sufficiently spread-out $J_\rho$,
together with the endpoint value, infinite right slope,
differentiability above the transition, and spread-out asymptotics.
The discussion following that conjecture also predicts the
fixed-interaction conclusions for every finite range.
Theorem~\ref{main:heights} states these conclusions for every admissible
$J$, including nearest neighbors, with the critical temperature
identified by Equation~\eqref{main:physical-threshold}.

The structural side of the argument also draws on a substantial
probabilistic literature.  Regev and Stephens-Davidowitz proved the
finite lattice Gaussian inequalities used below \citep{RSD2017}.
Aizenman, Harel, Peled, and Shapiro applied Gaussian monotonicity,
graph surgery, and level-set methods to depinning and the dual spin
phase \citep{AizenmanHarelPeledShapiro2022}.  Lammers and Ott developed
absolute-value correlation inequalities and delocalisation methods
\citep{LammersOtt2024}.  For a broader class of nearest-neighbor
square-lattice potentials, Lammers' dichotomy theory proves a logarithmic
variance lower bound with a universal positive coefficient throughout
the delocalized phase and shows that the delocalized set is closed in
the specified potential topology \citep{LammersDichotomy}.
The global depinning and dichotomy arguments in these works use planar
structure.  Our admissible finite-range graphs can be nonplanar, and
the structural proofs below cover them directly.
Recent critical results for hierarchical discrete Gaussian and related models
\citep{BiskupHuang2025} and Gaussian field limits for six-vertex heights
\citep{DuminilCopinKozlowskiLammersManolescu2026} concern different
interactions and further illustrate the range of height universality.

Villain introduced a periodic Gaussian edge weight as an approximation
to the planar cosine model.  The Villain duality and the analysis of
Jos\'e, Kadanoff, Kirkpatrick, and Nelson make the Gaussian and vortex
components explicit \citep{Villain1975,JoseKadanoffKirkpatrickNelson1977}.
The continuum field in our limit is imaginary Gaussian multiplicative
chaos.  Lacoin, Rhodes, and Vargas construct the purely imaginary case
for planar Gaussian free fields within their theory of complex chaos
\citep{LacoinRhodesVargas2015};
Junnila, Saksman, and Webb develop a general construction and moment
and regularity results \citep{JunnilaSaksmanWebb2020}.  The latter work
also proves that suitably normalized critical XOR-Ising spin smears with
plus boundary conditions on bounded simply connected planar domains
converge to smears of a conformally weighted real part of imaginary chaos.
Garban and Sep\'ulveda stated and sketched a joint limit of the
zero-boundary discrete GFF and its Wick exponential on a square, and
discussed the expected Villain connection
\citep[Proposition~6.3 and the discussion preceding Conjecture~3]{GarbanSepulveda2024}.
Theorem~\ref{main:spins} proves the full field statement for both
ordinary lattice spin models at sufficiently large fixed inverse
temperature.  This is distinct from the interface conjecture in that
work.  Newman and Wu's angular-gradient theorem uses a temperature
that tends to zero with the mesh \citep{NewmanWu2018}; here the
inverse temperature remains fixed.  We derive the smear moment estimates
from the finite-graph Lee--Yang theorems recalled by Newman and Wu,
including their detailed Villain proof, together with the scalar product
representation stated there \citep{NewmanWu2019}.

\subsection{Proof organization and the new estimates}

We first construct a structural height coefficient $a(k)$ with
\begin{equation}\label{main:coefficient-convention}
 k=\beta/v_J^2,\qquad \beff(J,\beta)=v_J^2a(k).
\end{equation}
Section~\ref{st:section} uses finite lattice Gaussian comparisons, free rectangular boxes, and folded
tests to yield Gaussian limits and the alternative $a=0$ or $a\ge8\pi$.
Section~\ref{an:section} uses an annular exploration to transfer the positive-coefficient limit to
Dirichlet domains, tori, and domains with pins, and to identify positivity
of $a$ with physical roughness.  The uniform estimates here must survive
exploration conditioning and pin events of small probability.

The local analytic calculus is developed in Section~\ref{rg:section}.
It keeps neutral, charged, boundary, and marked terms in distinct
estimates and supplies finite-range covariances and regulators with
explicit reserve.  Both spin constructions need local corrections that
remain controlled under repeated Gaussian integration, including near
inserted spins.  For Villain, Section~\ref{vi:section} uses exact duality
to initialize this expansion at the ordinary microscopic model.
For XY, the cosine interaction requires an additional estimate:
Section~\ref{xyh:section} represents it exactly and follows the
microscopic factors through successive integrations and local
subtractions.  The estimates distinguish a live Gaussian history from
completed scalar calculations, whose evaluated sizes multiply in later
operations.  Taylor remainders are
separated into fixed smooth directions and linear field functionals,
called atoms.  Estimates for these atoms control the additional
Gaussian sources.  A bound on the
evaluated history, together with separate charge estimates, supplies
sizes that compose through the local maps.  Section~\ref{xyf:section}
then chooses the three reference quadratic coefficients so that the
corrections decay.  An exact Hessian identity and a finite-dimensional
degree argument keep that choice inside the parameter region where
the history estimates hold.

In either spin model, marked estimates give all separated fundamental
spin correlations and a common insertion amplitude.  The full field law
then follows from a uniform two-point bound, Lee--Yang moment bounds,
collision estimates, and Sobolev tightness.  These steps retain the full
Green diagonal in Equation~\eqref{main:spin-limit}.

For heights near their physical endpoint, the weak-activity maps of
Section~\ref{wh:section} require physical initial data: one must first enter their domain from an arbitrary fixed
finite-range model.  Section~\ref{en:section} compares noisy block
averages of the original heights with Gaussian observations.  It gives
an exact local expansion of this density comparison, with activity as
small as the maps require whenever $a(k)>0$.  Mixed pinned domains and
complex Gaussian lattice comparisons supply control in the required
analytic norm.  Physical terminal observables select the admissible small
trajectory in Section~\ref{ep:section}.  Openness of the supermarginal branch identifies the
endpoint value; the marginal trajectory and its response then give
the infinite slope.  Figure~\ref{main:dependencies} records which
arguments must precede this last identification.

\begin{figure}[htbp]
\centering
\begin{tikzpicture}[
 box/.style={draw,rounded corners=2pt,align=center,text width=3.6cm,
             minimum height=1.45cm,inner sep=5pt,font=\small},
 arr/.style={-{Stealth[length=2mm]},line width=.5pt}]
 \node[box] (struct) at (0,0) {Finite comparisons\\structural coefficient\\
   \footnotesize Section~\ref{st:section}};
 \node[box] (annular) at (4.4,0) {Annular transfer\\and mixed pins\\
   \footnotesize Section~\ref{an:section}};
 \node[box] (entry) at (8.8,0) {Local height entry\\
   \footnotesize Section~\ref{en:section}};
 \node[box] (toolkit) at (0,-2.5) {Covariances, regulators\\local partition calculus\\
   \footnotesize Section~\ref{rg:section}};
 \node[box] (weak) at (4.4,-2.5) {Weak height maps\\physical diagnostics\\
   \footnotesize Sections~\ref{wh:section}, \ref{ep:section}};
 \node[box] (endpoint) at (8.8,-2.5) {Height endpoint\\and infinite slope\\
   \footnotesize Section~\ref{ep:section}};
 \node[box] (spin) at (0,-5) {Villain duality\\and XY histories\\
   \footnotesize Sections~\ref{vi:section}, \ref{xyh:section}};
 \node[box] (marks) at (4.4,-5) {Tuned flows\\marked correlations\\
   \footnotesize Sections~\ref{vi:section}, \ref{xyf:section}};
 \node[box] (field) at (8.8,-5) {Full spin-field laws\\and Villain bridge\\
   \footnotesize Sections~\ref{vi:section}, \ref{xyf:section}};
 \draw[arr] (struct)--(annular);
 \draw[arr] (annular)--(entry);
 \draw[arr] (entry)--(weak);
 \draw[arr] (toolkit)--(weak);
 \draw[arr] (weak)--(endpoint);
 \draw[arr] (toolkit)--(spin);
 \draw[arr] (spin)--(marks);
 \draw[arr] (marks)--(field);
 \draw[arr] (annular.south)--(weak.north);
\end{tikzpicture}
\caption{Dependencies among the main arguments.  The upper route
identifies the physical height endpoint; the lower route controls spin
insertions and then their full distribution.  Arrows show dependencies,
not reading order: the local entry construction is presented after the
weak maps whose parameters it accepts, and supplies their physical
initial data for the endpoint argument.  The Villain bridge also uses
the high-temperature height result of Section~\ref{wh:section}, and the
endpoint odd-sector argument returns directly to the structural face
ratio in Section~\ref{st:section}.}
\label{main:dependencies}
\end{figure}
